% META: {"id": "TEXP-ARI-EX-011", "chapitre": "TEXP-ARITHMETIQUE", "type_objet": "exercice", "capacites": ["TEXP-ARI-C2"], "capacites_codes": ["C2"], "methodes": ["M2"], "parcours": 2, "competences": ["raisonner"], "duree_min": 15, "mode_creation": "ex_nihilo", "sources_inspiration": [], "fichier_tex": "chapitres/TEXP-ARITHMETIQUE/exercices/TEXP-ARI-EX-011.tex", "corrige_tex": "chapitres/TEXP-ARITHMETIQUE/corriges/TEXP-ARI-CO-011.tex", "status": "generated"}
% BEGIN-VERIFY
% import math
% assert math.gcd(6, 15) == 3
% assert not any((6 * u) % 15 == 1 for u in range(15))
% for b in range(15):
%     a_des_solutions = any((6 * x) % 15 == b for x in range(15))
%     assert a_des_solutions == (b % 3 == 0)
% assert [x for x in range(15) if (6 * x) % 15 == 9] == [4, 9, 14]
% assert (2 * 3) % 5 == 1
% END-VERIFY
\begin{exercice}{TEXP-ARI-EX-011}{2}{15}
On considère la congruence $6x\equiv b\ [15]$.
\begin{enumerate}
  \item Calculer $\mathrm{PGCD}(6\,;\,15)$. Le nombre $6$ admet-il un inverse
        modulo $15$ ?
  \item Montrer que la congruence n'a de solution que si $3$ divise $b$.
  \item Résoudre $6x\equiv9\ [15]$ et préciser le nombre de solutions
        modulo $15$.
\end{enumerate}
\end{exercice}
